\documentclass[11pt,a4paper]{article}
\usepackage[left=22mm,right=22mm,top=20mm,bottom=21mm]{geometry}
\usepackage{amsmath,fontspec,unicode-math}
\setmainfont{lmroman10-regular.otf}[
  Path=fonts/, BoldFont=lmroman10-bold.otf,
  ItalicFont=lmroman10-italic.otf, BoldItalicFont=lmroman10-bolditalic.otf,
  Ligatures=TeX]
\setmonofont{lmmono10-regular.otf}[Path=fonts/]
\setmathfont{latinmodern-math.otf}[Path=fonts/]
\usepackage{booktabs,tabularx,array,graphicx,float,caption}
\usepackage{xcolor,fancyhdr}
\usepackage[numbers,sort&compress]{natbib}
\usepackage[colorlinks=true,linkcolor=navy,citecolor=navy,urlcolor=navy]{hyperref}
\definecolor{navy}{HTML}{13334F}
\definecolor{teal}{HTML}{00664F}
\definecolor{muted}{HTML}{47586A}
\definecolor{pale}{HTML}{EBF3F7}
\definecolor{ruleline}{HTML}{C9D9E2}
\hypersetup{pdftitle={Robust Sparse Smooth PCA for Closed-Set Face Recognition},pdfauthor={Paco, Sou Chak Tong}}
\newcommand{\doi}[1]{doi: \href{https://doi.org/#1}{\nolinkurl{#1}}}
\setlength{\parindent}{0pt}
\setlength{\parskip}{5pt}
\setlength{\headheight}{14pt}
\setlength{\fboxsep}{8pt}
\setlength{\emergencystretch}{2em}
\captionsetup{font=small,labelfont={bf,color=navy},justification=raggedright,singlelinecheck=false}
\graphicspath{{../results/faces/figures/}}
\pagestyle{fancy}
\fancyhf{}
\lhead{\small\color{muted}CISC3024 \enspace / \enspace AI Assignment \#3}
\rhead{\small\color{muted}November 2026}
\rfoot{\small\color{navy}\thepage}
\renewcommand{\headrulewidth}{0.3pt}
\newcommand{\reportfigure}[3]{%
  \begin{figure}[H]\centering
  \includegraphics[width=\linewidth]{#1}
  \caption{#2}\label{#3}
  \end{figure}}
\newcommand{\sourcequote}[3]{%
  \par\noindent{\small\texttt{#1}}\\
  \noindent\textit{\textquotedblleft#2\textquotedblright}\\
  \noindent #3\par}
\newcommand{\E}{\mathcal{E}}
\newcolumntype{Y}{>{\raggedright\arraybackslash}X}
\newcolumntype{P}[1]{>{\raggedright\arraybackslash}p{#1}}
% Generated by python -m src.face_summary.
\newcommand{\facemetrows}{%
Pixels & 256 & 86.25 & 76.50\,$\pm$\,4.37 & 63.25\,$\pm$\,4.20 & 58.09 \\
PCA & 16 & 82.50 & 68.25\,$\pm$\,5.05 & 55.50\,$\pm$\,5.42 & 50.91 \\
LDA & 16 & 83.75 & 81.50\,$\pm$\,6.58 & 68.25\,$\pm$\,8.46 & 64.98 \\
No smoothness & 16 & 82.50 & 68.00\,$\pm$\,4.01 & 56.50\,$\pm$\,4.37 & 52.03 \\
RSSPCA & 16 & 82.50 & 70.25\,$\pm$\,2.24 & 58.25\,$\pm$\,5.20 & 54.44 \\
}
\newcommand{\facemsrows}{%
PCA & 2.661 & 4.652 & 5.513 \\
No smoothness & 2.768 & 4.042 & 5.105 \\
RSSPCA & 2.845 & 3.513 & 4.162 \\
}
\newcommand{\facevalrows}{%
0.001 & 0.03 & 85.00 & 73.75 & 48.75 & 69.17 \\
0.001 & 0.3 & 85.00 & 77.50 & 51.25 & 71.25 \\
0.01 & 0.03 & 83.75 & 73.75 & 46.25 & 67.92 \\
0.01 & 0.3 & 86.25 & 76.25 & 51.25 & 71.25 \\
0.05 & 0.03 & 82.50 & 72.50 & 45.00 & 66.67 \\
0.05 & 0.3 & 85.00 & 76.25 & 47.50 & 69.58 \\
}

\begin{document}

{\bfseries\color{teal} CISC3024 --- Pattern Recognition \qquad AI Assignment \#3\par}
\vspace{2mm}
{\raggedright\bfseries\color{navy}
{\fontsize{29}{33}\selectfont Robust Sparse Smooth PCA\par}
{\fontsize{23}{28}\selectfont for Closed-Set Face Recognition\par}}
\vspace{2mm}
{\large Paco, Sou Chak Tong \enspace \textbf{U-C3-2652-3}}\par
{\small November 2026 \quad \textcolor{muted}{Technical report}}\par
\vspace{2mm}
{\color{ruleline}\hrule}
\vspace{3mm}
\textbf{Abstract.}
Robust sparse smooth PCA (RSSPCA) adds a pixel-grid penalty to an absolute-score criterion. On the 40-identity Olivetti database, we validate the 2025 update at $16\times16$ resolution, then compare projections on 80 held-out views of known identities. At 16 coordinates and 50\% corrupted training, RSSPCA reached 58.25\% accuracy versus 55.50\% for PCA, 56.50\% for its no-smoothing control and 68.25\% for Fisher LDA. Relative to PCA its clean-image projection error was 24.5\% lower in that condition. The original digit study provides secondary evidence about transfer.

\section{How AI Found the Algorithm}
The handout required AI assistance and six report topics. \texttt{AGENTS.md} defined the evidence and leakage rules. Prompt 01 governed the search and original Optdigits experiment. After reviewing that draft, the student asked for a face application to distinguish it from Assignment \#2, which also used $8\times8$ digits. Prompt 01B specifies the new feasibility gate and evaluation. Prompt 02 is the subsequent critical review.

\sourcequote{AGENTS.md}{Do not consult the final test set during this gate.}
{Codex downloaded only the development CSV. The independent arithmetic check and dataset smoke test ran before the selection record was written.}
\sourcequote{01\_RESEARCH\_IMPLEMENT\_REPORT.md}{do not begin the full implementation until the mandatory algorithm-selection gate (Section B2) is passed.}
{Codex screened primary papers, inspected the available reference solver and recorded a dated pass before creating the main pipeline.}
\reportfigure{face_workflow}{Instruction-to-execution workflow. Follow the top row from left to right and the lower row from right to left. The student made faces primary after the digit study. Prompt 02 reviewed both frozen records and regenerated the face report.}{fig:workflow}

\clearpage
\subsection{Search, screening and selection}
Codex ran a targeted search on 17 November 2026. The discovery criteria were a 2024--2026 original peer-reviewed method, a substantive PCA or LDA contribution, direct vision relevance and enough detail for a controlled Python evaluation. The initial web searches and solver follow-up used these exact strings.

{\small\setlength{\tabcolsep}{5pt}
\begin{tabularx}{\linewidth}{@{}P{21mm}Y@{}}
\toprule
Purpose & Query string \\
\midrule
RSSPCA & \texttt{2025 robust sparse smooth principal component analysis RSSPCA face recognition} \\
Sparse LDA & \texttt{2024 robust and sparse linear discriminant analysis face image objective robust sparse LDA} \\
2D methods & \texttt{2025 2026 robust two dimensional principal component analysis face recognition algorithm sparse Lp} \\
Broader search & \texttt{2024 2025 novel two dimensional PCA image recognition generalized weighted 2DPCA face recognition open access algorithm} \\
LDA full text & \texttt{"Towards robust and sparse linear discriminant analysis for image classification" pdf RSLDA+ Liu Feng} \\
Code & \texttt{DLR-2DPCA Pengfei Bi GitHub double layer robust code} \\
\bottomrule
\end{tabularx}}\par
A literature lookup recovered the SCRAMBLE preprint. Codex used the journal article to verify the method and wear-scar image application. The complete query ledger is in \path{docs/research.md}. The literature-review skill guided primary-source screening. The computation-audit skill guided the independent calculation. The experimental-design and statistical-analysis skills informed the shared corruption masks and the interpretation of variation across runs \citep{kassis2026}.

\sourcequote{02\_CRITICAL\_REVIEW\_AND\_FIGURES.md}{Preserve factual integrity and valuable negative findings.}
{Codex rechecked the frozen predictions, repaired two prose values and plotted the archived stopping traces. The held-out results, including the weaker results against LDA and on digits, were retained.}

\begin{table}[H]\small
\caption{Shortlist and internal suitability scores. Weights for recency, PCA/LDA centrality, implementation, experimental fairness, mathematical verification and figure potential were 15/20/20/20/15/10 percent. These are project decision weights.}
\begin{tabularx}{\linewidth}{@{}P{26mm}YP{56mm}r@{}}
\toprule
Candidate & Contribution and vision use & Decision evidence & /5\\
\midrule
RSSPCA \citep{wang2025} & Absolute projected dispersion, sparse and grid-smooth loadings. Face recognition and reconstruction. & Open equations and MATLAB reference. Direct smoothing ablation at fixed feature count. & 5.00\\
DLR-2DPCA \citep{bi2025} & Cut reconstruction loss and adaptive weights. Face, digit and traffic sign recognition. & Readable algorithm. More numerical boundary cases. Matrix outputs restrict matched budgets. & 4.30\\
SCRAMBLE \citep{pfeiffer2025} & Robust cellwise reconstruction and elastic net. Wear-scar image regression. & Open R reference. More solver and scale choices for transfer to digit recognition. & 4.15\\
RSLDA+ \citep{liu2024} & Sparse LDA with ADMM and thresholding. Image classification. & Complete solver unavailable through sources accessed during the search. & 2.95\\
\bottomrule
\end{tabularx}
\end{table}
\textbf{Decision gate.}
Wang et al. published RSSPCA in PLOS One in May 2025. Its four-neighbor graph gives a concrete mechanism to test on raster data. At a fixed coordinate budget, the graph penalty can be switched off while keeping the absolute-score objective, sparsity setting and classifier. This made it possible to compare the effect of spatial regularization separately from the overall comparison with PCA and LDA.

\clearpage
\subsection{Why RSSPCA passed the gate}

Recency and a direct change to PCA or LDA were necessary conditions. For the remaining candidates, the question was whether the published equations could support an independently checked implementation and a fair experiment. Table~\ref{tab:gate-scores} shows the individual ratings behind the shortlist.

\begin{table}[H]\small\centering
\caption{Recorded gate ratings from 1 to 5. A weighted score averages the six ratings with the indicated weights. Eligibility also requires an accessible solver.}
\label{tab:gate-scores}
\setlength{\tabcolsep}{5pt}
\begin{tabular}{@{}lrrrrrrr@{}}
\toprule
Candidate & Recency & Centrality & Feasibility & Fairness & Math & Figures & Score\\
Weight (\%) & 15 & 20 & 20 & 20 & 15 & 10 & \\
\midrule
RSSPCA & 5 & 5 & 5 & 5 & 5 & 5 & 5.00\\
DLR-2DPCA & 5 & 5 & 3 & 4 & 5 & 4 & 4.30\\
SCRAMBLE & 5 & 5 & 3 & 4 & 4 & 4 & 4.15\\
RSLDA+ & 5 & 5 & 1 & 2 & 2 & 3 & 2.95\\
\bottomrule
\end{tabular}
\end{table}

DLR-2DPCA was a viable alternative with a readable algorithm and direct digit-recognition precedent. On the $8\times8$ grid, a right projection with $r$ columns would deliver $8r$ scalar scores, which restricts the dimensions that can be matched to vector PCA and LDA. SCRAMBLE also had an accessible method and reference package. Its wear-scar image regression used cellwise robustness with sparsity disabled. Transferring the full method here would require choices about the cellwise loss, scale and sparsity fitting. The original search had only the RSLDA+ publisher preview. The final audit located an \href{https://github.com/EMXlight/RSLDAplus}{author-linked MATLAB solver}. The recorded scores describe the evidence available when the selection was locked.

\subsection{What the development-only probe measured}

The two-pixel calculation independently solved the defining linear system by Cramer's rule. The numerical residual was $6.28\times10^{-16}$ and the smallest eigenvalue of the symmetric system was $0.873$. A separate calculation with an exactly zero loading coordinate checked the branch most exposed to the inverse-diagonal formulation. The normalized update and its stopping interpretation are derived in Section~\ref{sec:method}.

The data check read only the 3,823-row development file. It found 3,823 unique image vectors, ten classes with 376 to 389 examples each, and two constant coordinates. On the 3,058/765 feasibility split, PCA with eight coordinates and Euclidean 3-NN reached 96.08\%. A one-direction RSSPCA probe met the $10^{-4}$ step threshold in 27 iterations. These checks established that the representation, evaluator and update could run on the available CPU before the comparison study was designed. The 96.08\% value refers to that feasibility split.

The original gate and run records date the method decision before implementation in \path{docs/algorithm_selection.md}. The separate \path{results/selection.json} records the dimension and selected penalties. The execution chronology comes from the staged runner and project records. The preserved Git snapshots were created after their respective runs.

During Prompt 02 the academic-review skill guided examiner-style checks of the equations, experimental claims and figures. The scientific-visualization and Nature Figure skills informed the matched-mask contrasts and the stopping plot. The English LaTeX and PDF guidance supported consistent typography and page-level inspection. This was a Codex self-audit with additional AI review passes.

\clearpage
\subsection{Face feasibility before changing the primary experiment}

The pivot retained the 2025 method and moved the primary question from digit classification to closed-set face identification. The \texttt{study/optdigits-baseline} branch records the verified digit project. The face branch adds a separate data loader and runner, preserving the original solver and digit artifacts. The dataset choice uses the accessible $64\times64$ derivative of the AT\&T Laboratories Cambridge archive \citep{attcam,sklearnfaces}.

\sourcequote{01B\_FACE\_PRIMARY\_MIGRATION.md}{Preserve the mathematically validated RSSPCA mechanism.}
{The face runner reuses the validated update and deflation. Its training-only gate, split, tuning and results are recorded separately.}

At $d=4{,}096$ the per-update dense Cholesky term scales by $16^3$ relative to $d=256$. A single float64 $d\times d$ matrix needs 128 MiB instead of 0.5 MiB. The gate therefore tested the prespecified $16\times16$ representation produced by block averaging. The shared implementation accepts a rectangular shape and builds a four-neighbor graph in the same row-major order as the input.

The gate split only the intended fitting partition into 160 fit and 80 internal check images, with every identity present in both. It compared edge energy with a direct sum of vertical and horizontal squared differences. It checked zero-coordinate preservation and sign invariance, and compared the symmetric and independent nonsymmetric updates. The normalized difference was $3.25\times10^{-16}$ and the relative symmetric-system residual was $3.85\times10^{-16}$. The rank-24 fit reproduced the loading array exactly, with residual overlap $\|R W\|_{\mathrm F}=3.18\times10^{-14}$. The test also checked projector idempotence and that centered projection improves or matches reconstruction from the mean alone.

\begin{table}[H]\small\centering
\caption{Face gate measurements on the fitting partition only. The final validation and test partitions were excluded. Times are single-thread projection fits.}
\begin{tabular}{@{}lrrr@{}}
\toprule
Method & Coordinates & Fit time (ms) & Internal accuracy (\%) \\
\midrule
Raw pixels & 256 & 1.0 & 73.75\\
PCA & 24 & 5.0 & 67.50\\
Fisher LDA & 24 & 5.3 & 93.75\\
No smoothness & 24 & 454.7 & 70.00\\
RSSPCA & 24 & 391.1 & 68.75\\
\bottomrule
\end{tabular}
\end{table}

The fitting matrix had numerical rank 159 and occupied 320 KiB in float64. The repeated 24-coordinate RSSPCA fit peaked at 2.84 MiB of allocations tracked by Python's tracer. That measurement excludes native scratch space outside the tracer. The observed fit stayed below the gate bounds of 10 seconds and 64 MiB, while the matched PCA and LDA arms exceeded 35\% internal accuracy. The locked protocol in \path{docs/face_study_design.md} governed the subsequent evaluation.

\reportfigure{face_mechanism}{Face pipeline validated by the gate. Iteration uses the current residual and the pixel graph in the symmetric positive-definite (SPD) solve. The accepted loading expands the deflated span. Held-out readouts use the saved mean and raw scores for 3-NN or the orthogonal span for reconstruction.}{fig:solver}

\clearpage
\section{Algorithm Description}
\label{sec:method}
\subsection{Classical reference subspaces}

Let $a_i\in\mathbb R^d$ be training image $i$, flattened in row-major order, with identity $y_i\in\{0,\dots,C-1\}$. There are $n$ training images, $d=256$ spatial coordinates and $C=40$ classes in the face study. The training mean is $\mu=n^{-1}\sum_i a_i$. Let $X\in\mathbb R^{n\times d}$ have rows $x_i^\top=(a_i-\mu)^\top$.

Classical PCA \citep{pearson1901} solves $\max_{\|w\|_2=1} w^\top S w$ for $S=X^\top X/(n-1)$. Differentiating with the norm constraint gives $Sw=\lambda w$. With thin SVD $X=U\Sigma V^\top$, the first $k$ eigenvectors are the first $k$ columns of $V$, with eigenvalues $\sigma_j^2/(n-1)$ for decreasing singular values $\sigma_j$. The rank is at most $\min(d,n-1)$, and the experiment uses unwhitened scores. Equal eigenvalues identify an eigenspace, with freedom to rotate the basis inside it. Isolated eigenvectors have an arbitrary sign.

For Fisher LDA \citep{fisher1936}, let $\mu_c$ and $n_c$ be the class mean and size. The within- and between-class scatters are
\[
 S_W=\sum_{i=1}^{n}(a_i-\mu_{y_i})(a_i-\mu_{y_i})^\top,\qquad
 S_B=\sum_{c=0}^{C-1}n_c(\mu_c-\mu)(\mu_c-\mu)^\top.
\]
When $S_W$ is positive definite, maximizing $w^\top S_B w$ under $w^\top S_Ww=1$ yields $S_Bw=\lambda S_Ww$. The weighted class means sum to zero after global centering, so $\operatorname{rank}(S_B)\le C-1=39$. Also $\operatorname{rank}(S_W)\le n-C$, so the 240-image tuning fit has at most 200 within-class directions in 256 dimensions. The rank-aware SVD fit whitens retained within-class coordinates before extracting the directions of class-mean separation. This avoids forming $S_W^{-1}$ and defines the baseline's score geometry.

\subsection{RSSPCA objective and the image graph}

For extraction step $j$, let $R\in\mathbb R^{n\times d}$ be the residual training matrix, with $R=X$ initially. Wang et al.\ \citep{wang2025} define the loading objective
\begin{equation}
 \max_w \|Rw\|_1
 \quad\text{subject to}\quad
 \|w\|_2=1,\qquad \|w\|_1\le c_1,\qquad w^\top Lw\le c_2 ,
 \label{eq:objective}
\end{equation}
where $c_1,c_2$ are sparsity and roughness bounds. Absolute projected scores grow linearly with projection magnitude, compared with quadratic growth in PCA. This follows the earlier PCA-$\ell_1$ criterion \citep{kwak2008}. For example, doubling one residual row at fixed $w$ doubles its absolute-score contribution and quadruples its squared-score contribution. A large residual may still affect the direction because the sign vote multiplies the residual row.

Build $A$ by joining horizontal and vertical neighbors and use $L=\operatorname{diag}(A\mathbf1)-A$. A $h\times b$ grid has $h(b-1)+(h-1)b$ undirected edges, 480 for the $16\times16$ grid. There is no wrap around at a row boundary. For edge set $\E$,
\begin{equation}
 w^\top Lw=\tfrac{1}{2}\sum_{p,q}A_{pq}(w_p-w_q)^2
           =\sum_{\{p,q\}\in\E}(w_p-w_q)^2 .
 \label{eq:energy}
\end{equation}
Thus $L$ is symmetric positive semidefinite and annihilates constant vectors. Penalizing this energy favors nearby pixel weights that agree. The $\ell_1$ penalty on unit vectors favors concentration on fewer pixels. These effects can compete along image boundaries where useful coefficients change quickly.

The graph in Figure~\ref{fig:solver} is the same for every image and training condition. The stabilized solve below weights the coupling by the current coefficient magnitudes, so its effective strength decreases as either coefficient approaches zero. The no-smoothing arm measures what changes when the graph term is removed.

\clearpage
\subsection{From the minorizer to the published update}

For a current unit direction $w$ with nonzero coordinates, set $D=\operatorname{diag}(|w|)$ and $v=R^\top\operatorname{sign}(Rw)\in\mathbb R^d$, using $\operatorname{sign}(0)=0$. For any trial $z\in\mathbb R^d$,
\[
 \|Rz\|_1\ge v^\top z,\qquad
 |z_p|\le\frac{z_p^2}{2|w_p|}+\frac{|w_p|}{2}.
\]
Both inequalities are tight at $z=w$. The first follows termwise from $|b|\ge s b$ for $|s|\le1$. The second follows by expanding $(|z_p|-|w_p|)^2\ge0$.

Write the Lagrangian with penalty multipliers $a,b\ge0$ on the last two constraints in Eq.~\eqref{eq:objective} and $c>0$ on $\|z\|_2^2-1$. Replacing the two absolute-value terms with their bounds gives the concave quadratic minorizer, up to an additive constant,
\[
 G(z\mid w) = v^\top z -\tfrac a2 z^\top D^{-1}z -b z^\top Lz-c z^\top z.
\]
Its unconstrained stationary point solves $(2cI+aD^{-1}+2bL)z=v$. The published update normalizes the direction, absorbing the positive scale $2c$ by defining $\eta=a/(2c)$ and $\gamma=b/c$. The pre-normalization vector $p$ solves
\begin{equation}
 (I+\eta D^{-1}+\gamma L)p=v,\qquad
 w^+=p/\|p\|_2 .
 \label{eq:original}
\end{equation}
For invertible $D$, putting $p=Du$ gives the published zero-safe form $p=D(D+\eta I+\gamma LD)^{-1}v$. Our implementation uses $H=D^{1/2}$ and the equivalent symmetric system
\begin{equation}
 M t=Hv,\quad M=D+\eta I+\gamma H L H,\qquad
 p=Ht .
 \label{eq:spd}
\end{equation}
Set $B=D+\eta I+\gamma LD$. When some entries of $D$ vanish, order the nonzero coordinates first. Then $B$ is block lower triangular because each zero coordinate contributes only $\eta$ on its column diagonal. The nonzero block is similar to the corresponding positive-definite block of $M$ via $H$ restricted to those coordinates. The other diagonal block is $\eta I$, so $B$ is invertible. Now $HB=MH$ gives $HB^{-1}v=M^{-1}Hv$ and both formulations yield the same $p$. Since $M\succeq\eta I$, $\|L\|_2\le8$ by the maximum grid degree four, and $\|D\|_2\le1$, the condition number obeys $\kappa_2(M)\le(\eta+1+8\gamma)/\eta$. This is 3401 at the chosen face setting. The code uses a Cholesky-based solve.

\subsection{Worked example and stopping interpretation}

Take rows $(1,2)$ and $(-1,-2)$, $w=(0.6,0.8)^\top$, and a two-pixel edge with $L=\left[\begin{smallmatrix}1&-1\\-1&1\end{smallmatrix}\right]$. Then $v=(2,4)^\top$. For $\eta=0.2$ and $\gamma=0.5$, the nonsymmetric system for $u$ has rows $(1.1,-0.4)$ and $(-0.3,1.4)$. Cramer's rule followed by $p=Du$ gives
\[
 p=(132/71,\,200/71)^\top,\qquad w^+=(0.550842,\,0.834609)^\top.
\]
The SPD computation agreed. Starting with $(1,0)^\top$ gave $p=(20/17,0)^\top$, illustrating that exact zero coordinates are absorbing for $\eta>0$.

The method maximizes the unconstrained quadratic before normalization. A sphere-constrained quadratic optimizer would also satisfy tangent stationarity. The numerical example gives $\|(I-w^+{w^+}^\top)(v-Kw^+)\|_2=0.376$ for $K=I+\eta D^{-1}+\gamma L$, so the normalized candidate fails that stationarity condition. Our implementation follows the published fixed-point map and selects fixed penalty ratios $\eta$ and $\gamma$. It does not estimate multipliers to enforce the explicit budgets $c_1,c_2$. The operational stop checks the size of the latest update. Section~\ref{sec:face-diagnostics} tests what this stop signifies in the actual fits.


\clearpage
\section{How AI Implemented the Algorithm}

Codex inspected the authors' GPL-3.0 MATLAB reference \citep{matlab} and implemented the paper equations independently with NumPy and SciPy. The reference initializes with axes from the original SVD, stops using a relative change in the L2 score norm and uses accumulated raw directions in deflation. This project initializes each component with the leading covariance eigenvector of the \emph{current residual}. It uses the sign-invariant step size $\min(\|w^+-w\|_2,\|w^++w\|_2)$ and stops at $10^{-4}$ or 100 updates. The code records both step size and mean absolute score. The largest-magnitude entry is made nonnegative for display. With equal leading eigenvalues, the initialization can depend on the numerical basis picked by the solver.

\subsection{Deflation, coordinate scores and cost}

Write $W_j=[w_1,\ldots,w_j]$ and let $Q_j$ have orthonormal columns spanning $W_j$, maintained by twice removing the existing basis projection from each new column before normalizing. We recompute \(R_j=X-(XQ_j)Q_j^\top\). Then $R_j W_j=0$ up to roundoff. The expression $I-W_jW_j^\top$ is an orthogonal projector only when the columns of $W_j$ are orthonormal. The raw columns can remain oblique, so prediction uses $W_k^\top(a-\mu)$ and retains the learned metric on those coordinates. The implementation rejects excess dimensions, exhausted residuals, dependent new directions and an update with zero norm.

For two images $a$ and $b$, squared Euclidean distance between the raw scores is
\[
 \|W_k^\top(a-b)\|_2^2=(a-b)^\top W_kW_k^\top(a-b).
\]
The common mean cancels in this calculation. Flipping the sign of a loading leaves this distance unchanged. Using the orthonormal basis $Q_k$ for classification would instead assign equal weight to all directions in the learned span and could change which neighbors vote. The two bases therefore have distinct roles in the implementation.

For $k$ components and up to $T$ updates per component, the dense implementation costs $O\!\left(k(nd^2+d^3)+kT(nd+d^3)+ndk^2\right)$ time. The terms bound residual covariance/eigenvector initialization, sign votes and SPD solves, and growing orthogonal projections. Working storage is $O(nd+d^2+dk)$ plus diagnostic traces. A test score costs $O(dk)$ per image and brute-force 3-NN compares $k$-coordinate scores against $n$ training examples. The primary face pipeline uses $256\times256$ systems, measured in the training-only gate.

\begin{table}[H]\small
\caption{Pseudocode and equation-to-code mapping. Fit inputs are training rows, labels, penalties and the feature budget. Each returned face loading has 256 entries.}
\begin{tabularx}{\linewidth}{@{}P{9mm}YP{53mm}@{}}
\toprule
Step & Operation & Implementation \\
\midrule
1 & Compute $\mu$, center training rows, set $R\gets X$, and form $L$. & \texttt{fit\_rsspca}, \texttt{grid\_laplacian} \\
2 & For each component, initialize $w$ from the leading residual covariance eigenvector. & \texttt{fit\_rsspca} \\
3 & Form $v,D,H,M$, solve Eq.~\eqref{eq:spd}, normalize, and repeat to the step threshold or cap. & \texttt{fixed\_point}, \texttt{iterate\_component} \\
4 & Append $w$ to $W$ and reproject $X$ away from its span. & \texttt{fit\_rsspca} \\
5 & Project train and test with the saved $\mu,W_k$. Vote among the 3 closest training scores. & \texttt{Projection.transform}, \texttt{classifier} \\
6 & For reconstruction, compute a thin QR basis for the span of $W_k$ and project the clean input. & \texttt{reconstruct} \\
\bottomrule
\end{tabularx}
\end{table}

\clearpage
\subsection{Reconstruction in a nonorthogonal loading span}

For a face column vector $a$, write $W_k=Q_kB_k$ using thin QR. With linearly independent columns in $W_k$, $B_k$ is invertible, so $Q_k$ and $W_k$ span the same subspace. Minimizing $\|(a-\mu)-Q_k z\|_2^2$ gives $z=Q_k^\top(a-\mu)$ by the normal equations. Thus
\[
\widehat a=\mu+Q_kQ_k^\top(a-\mu),\qquad
\mathrm{MSE}(a,\widehat a)=d^{-1}\|a-\widehat a\|_2^2 .
\]
The residual is orthogonal to the span. The projector $Q_kQ_k^\top$ is symmetric and idempotent. A small test with deliberately nonorthogonal columns checked reconstruction against an independent least-squares solve and verified that a second projection leaves the reconstruction unchanged. The test target is the same clean reduced image supplied to the projection. Noise acts only on the training rows used to learn $\mu$ and $W_k$.

\subsection{Verification and execution}

The combined suite passed 13 tests. It includes the ten Prompt 01 regression checks and three face additions for deterministic identity splits, rectangular box averaging and the nonorthogonal reconstruction check. The training-only face probe independently tested the larger grid and all recognition arms.

\begin{table}[H]\small
\caption{Representative verification checks from \texttt{tests/test\_pipeline.py}. Numerical tolerances are absolute floors used by the assertions.}
\label{tab:checks}
\begin{tabularx}{\linewidth}{@{}P{33mm}Y@{}}
\toprule
Check & Evidence and the failure it can expose\\
\midrule
Graph geometry & A $2\times3$ grid has seven unique edges, zero row sums, a nonnegative spectrum up to roundoff and no edge spanning a row boundary. A hand-summed edge energy checks the flattening convention.\\
Update equivalence & The SPD solve is compared with an independent nonsymmetric solve at three penalty pairs, including $\eta=10^{-9}$ with $\gamma=0$ and $\eta=10^{-8}$ with $\gamma=100$. The test uses a $10^{-10}$ absolute tolerance floor after aligning signs and verifies that exact zero coordinates remain zero.\\
Deflation and prefix & Tests compare repeated loading fits and the first loading extracted on its own. An independently computed basis gives a residual orthogonal to the learned span within $10^{-13}$. A rank-deficient input must raise.\\
Classical baselines & PCA is checked against the scikit-learn subspace and the covariance eigen-equation. The Fisher SVD subspace is checked against a generalized eigenproblem and adding a constant pixel must preserve its scores up to sign.\\
Data and selection & The two corruption severities share row IDs, patch locations and values on their intersection. A serialization round trip checks mask replay. A tie in mean accuracy exercises the stated selection order.\\
\bottomrule
\end{tabularx}
\end{table}

\textbf{Replay and metric checks.}
The face run used Windows 11, Python 3.14.3, NumPy 2.4.4, SciPy 1.18.0 and scikit-learn 1.9.0 on CPU, with BLAS and OpenMP limited to one thread. A new replay matched the original settings, split, loadings, final masks, 143 prediction arrays and 99 reconstruction-error arrays exactly. Fit times are measured wall times and can vary.

The Prompt 02 audit independently decoded the source images, regenerated both mask streams and computed Euclidean distances and class votes directly. It reproduced all 11,440 classifications and used least squares to verify the per-image reconstruction errors. Across all fits, the normalized PCA eigen-equation residual stayed below $2.5\times10^{-15}$, and the retained LDA scores satisfied the within-class whitening check. The executable checks and diagnostic record are in \path{scripts/review_face.py} and \path{results/faces/prompt02_audit.json}.

\clearpage
\section{Experiment Settings and Results}
\subsection{Faces, split and training corruption}

The primary dataset is the $64\times64$ Olivetti derivative of the AT\&T Laboratories Cambridge database \citep{attcam,sklearnfaces}. It has ten grayscale images for each of 40 identities. Acquisition varies in expression, lighting and details such as glasses. The released images are cropped frontal views against a largely uniform background. This experiment evaluates closed-set recognition, so fitting, validation and test all contain every identity. Its question is whether a new image of a known person can be matched to that identity. The experiment does not evaluate unseen identities.

The input orientation follows the scikit-learn loader. Each source byte is divided by 255, and nonoverlapping $4\times4$ source blocks are averaged to produce one $16\times16$ image. This is a fixed box filter followed by integer downsampling. It attenuates high frequencies before decimation and does not use any face-test statistics. All arms use the same reduced pixels and the mean fitted on their own training rows.

\reportfigure{face_protocol}{The first four archived perturbations in the seed-8301 severe mask applied to the final fitting pool. Pixel positions are zero indexed on the $16\times16$ grid. Colors distinguish untouched cells and cells assigned either zero or one. A replacement can equal the original value.}{fig:protocol}

With split seed 20261117, six images per identity go to fitting, two to validation and two to test. The 240/80/80 source rows are disjoint and the final audit found 400 distinct reduced images. Tuning used clean validation images and 240 fitting examples. Final fits used the union of fitting and validation rows, eight images per identity. The clean test denominator is 80 throughout.

A mask selects $\lfloor fn\rfloor$ training rows without replacement for $f\in\{0,.25,.50\}$. Each selected $16\times16$ face receives an $8\times8$ patch whose location is drawn uniformly among 81 positions. Within it, each value is independently set to zero or one with equal probability. With 320 final training rows, 80 or 160 images receive patches. The lower severity is a prefix of the higher severity mask within each seed. All arms share the row identities, patch values, locations and clean test images.

\clearpage
\subsection{Matched comparison, tuning and outcome definitions}

The four reduced arms use 8, 16 and 24 coordinates, with 16 the primary budget. Unprojected pixels use all 256 input coordinates. All models share the same training rows, masks and uniform Euclidean 3-NN. Fisher LDA uses identity labels and its fitted within-class scaling. The PCA-family arms fit without labels. All arms use labels at the voting stage and retain their fitted coordinate scales. A three-way class-vote tie goes to the smallest encoded identity. The review found no tied distances at the third-neighbor boundary in the saved evaluations.

The development grid uses three sparsity settings $\eta$ and two graph settings $\gamma$. A candidate uses one shared training mask stream per severity, seed 8101. The selected pair maximizes the mean of the three accuracies at 16 coordinates, with ties resolved by smaller $\gamma$ then smaller $\eta$. In Table~\ref{tab:faceval}, the two leading settings each classify 171 of the 240 repeated validation predictions correctly. The rule selected $(0.001,0.3)$. The controlled ablation uses $\eta=0.001$ and $\gamma=0$.

\begin{table}[H]\small\centering
\caption{Accuracy (\%) on the 80 development validation rows. The three columns vary only the training corruption.}
\label{tab:faceval}
\begin{tabular}{@{}rrrrrr@{}}
\toprule
$\eta$ & $\gamma$ & Clean & 25\% & 50\% & Mean \\
\midrule
\facevalrows
\bottomrule
\end{tabular}
\end{table}

At $\gamma=0$ the system in Eq.~\eqref{eq:spd} is diagonal, with $p_i=|w_i|v_i/(|w_i|+\eta)$. This makes the attenuation of small coefficients explicit. Positive $\gamma$ adds a degree-dependent diagonal penalty and couples neighboring entries through $-\gamma\sqrt{|w_iw_j|}$. Removing the graph term at fixed $\eta$ isolates its contribution under the chosen settings. The separate effects of the $\ell_1$ score and sparsity term are outside this ablation.

Let $C_{ij}$ be the test confusion count for true identity $i$ and predicted identity $j$. Accuracy is $\sum_i C_{ii}/80$. Balanced accuracy averages $C_{ii}/\sum_j C_{ij}$ over the 40 identities and equals accuracy here because every identity has exactly two views. Macro-F1 averages $2C_{ii}/(\sum_j C_{ij}+\sum_j C_{ji})$. The latter also reflects whether a classifier overpredicts a label. In the secondary reconstruction analysis, we average pixel MSE across the same 80 clean views for each PCA-family fit.

The five final training masks use seeds 8301--8305 at each positive severity. Clean training is deterministic and runs once. Before test scoring, the face runner saved and hashed the selected configuration, split, fit records and matrices, as well as recording source code and data digests. The test phase checks that freeze and reconstructs the training exemplars from the archived mask arrays. In the recorded run, the freeze completed at 14:59:22 UTC+8 on 17 November 2026 and final scores were written at 14:59:40.

\clearpage
\subsection{Held-out identification}

\begin{table}[H]\small\centering
\caption{Locked study holdout accuracy (\%). Positive-severity entries are mean $\pm$ sample SD over five training masks. Macro-F1 is averaged at 50\% corruption.}
\label{tab:facemet}
\begin{tabular}{@{}lrrrrr@{}}
\toprule
Method & Coordinates & Clean & 25\% & 50\% & Macro-F1 at 50\%\\
\midrule
\facemetrows
\bottomrule
\end{tabular}
\end{table}

\reportfigure{face_recognition}{Identification on the same 80 clean test faces. A, clean accuracy and the means (lines) and individual masks (small points) at positive severities. Raw pixels use 256 coordinates and other arms use 16. B, paired accuracy differences in percentage points for five masks per severity. Open symbols mark mean differences. Positive values favor RSSPCA.}{fig:recognition}

The three PCA-family arms each classified 66 of 80 clean views correctly. At 25\% corruption, RSSPCA improved over PCA by 2.00 percentage points on average (paired range $-3.75$ to $+6.25$). At 50\% the difference was $+2.75$ points (range $-7.50$ to $+10.00$). In the severe condition it beat PCA on three masks, tied on one and lost on one. This supports an average advantage under these masks while showing that the sign of the effect is sensitive to the corruption realization.

Adding smoothing at fixed $\eta$ gave $+2.25$ points at 25\% and $+1.75$ at 50\%. The respective paired ranges were $-1.25$ to $+5.00$ and $-5.00$ to $+6.25$. At 50\%, Fisher LDA led RSSPCA by 10.00 points on average. Raw pixels also scored higher than RSSPCA but preserved sixteen times as many coordinates. The supervised direction criterion and larger raw representation change what information is retained, even though the neighbors and training masks are held fixed.

One additional correct face changes accuracy by 1.25 percentage points. The agreement of conclusions across masks is more informative here than a single comparison. The sample SD describes mask sensitivity on these fixed images; it gives no interval for an independent sample of people or test views.

\clearpage
\subsection{Spatial loadings and the smoothing ablation}

\reportfigure{face_loadings}{First four unit loading maps in each method's extraction order, from the 50\% training mask at seed 8301. Loadings shown in the same column are not matched across methods. The largest absolute coefficient is oriented nonnegative. All maps share a zero-centered weight scale. Row annotations average grid energy over the first 16 directions.}{fig:faceload}

At 50\%, RSSPCA averaged grid energy 1.226 versus 2.440 without smoothing. For unit directions this is the sum of squared coefficient differences on the 480 edges. It dropped about 50\% when the graph term was enabled. The fraction of coefficients with absolute value at most $10^{-3}$ fell from 2.67\% to 1.97\%, averaging over the five masks at 16 coordinates. The chosen threshold counts small weights, so it describes coefficient magnitudes rather than exact support recovery. These changes are consistent with smoothing spreading weight across adjacent pixels. The paired recognition contrasts in Figure~\ref{fig:recognition} measure the smaller change in identification.

The first loading of each arm shows a broad region with a common sign. Later loadings have more localized sign changes. At an edge in a map, the penalty increases with the squared difference between neighboring weights, even if both weights are small. The single-seed maps illustrate where such differences occur. The distribution of edge energies across all five severe masks in Figure~\ref{fig:facediag} tests whether the reduction holds beyond that illustration.

\clearpage
\subsection{Fitting cost and stopping behavior}
\label{sec:face-diagnostics}

\reportfigure{face_diagnostics}{Diagnostics at 50\% corrupted training. A, mean edge energy at 16 coordinates for each of five masks. Outlined markers give means. B, measured fit time on a logarithmic axis with median markers. Pixels use 256 coordinates, other arms 16. C, empirical cumulative percentage of loadings stopped after each number of updates. Each robust arm contributes 24 loadings from each of five fits. The step threshold is $10^{-4}$ and the cap of 100 is beyond the plotted range.}{fig:facediag}

Both robust arms used the direction-change stop for all 264 extracted directions across the complete study. In the severe fits the median was 10.5 updates for RSSPCA and 10 for the no-smoothing arm. The review independently took one more update from each saved direction. For RSSPCA, one subsequent step slightly exceeded $10^{-4}$, with a maximum of $1.0235\times10^{-4}$. Also, 245 of its 264 traces contained an absolute-score decrease of more than $10^{-12}$ between successive iterates. The logged score excludes the penalties. These diagnostics quantify finite-step termination and leave objective convergence and global optimality unresolved.

Median fit time at 16 coordinates was 153.0 ms for RSSPCA, 142.3 ms without smoothing, 9.2 ms for PCA and 11.7 ms for LDA. Timing includes projection fitting and preparing the 3-NN training scores. The PCA and LDA decompositions were performed once per mask and charged at each prefix, while the robust arms were charged their cumulative extraction time through each prefix. Timings characterize this implementation and CPU. Prediction medians were 0.35--0.38 ms for the 80-view batch across those four arms.

\clearpage
\subsection{Feature budget and reconstruction fidelity}

\reportfigure{face_budget}{Sensitivity to coordinate budget at 50\% corrupted training. The dashed guide marks the primary budget of 16. Points and whiskers give means and sample SDs over five masks on the same 80 clean test views. A, identification. B, mean per-pixel reconstruction MSE multiplied by $10^3$, measured before clipping.}{fig:facebudget}

At 50\%, increasing the RSSPCA budget from 8 to 16 to 24 raised accuracy from 57.00\% to 58.25\% to 59.75\%. LDA improved from 59.75\% to 68.25\% to 70.75\%. The comparison at 24 coordinates is a prespecified sensitivity result. The primary budget remains 16. Increasing the budget lets additional variation reach the classifier, and the different gains show that useful identity information is distributed differently across the fitted directions.

\begin{table}[H]\small\centering
\caption{Mean clean-image projection MSE $\times10^3$ at 16 coordinates. Image scale is $[0,1]$, and errors are measured before clipping. Each nonzero severity averages five fits on the same 80 test images.}
\label{tab:facemse}
\begin{tabular}{@{}lrrr@{}}
\toprule
Method & Clean training & 25\% & 50\%\\
\midrule
\facemsrows
\bottomrule
\end{tabular}
\end{table}

With clean training, PCA reconstructed best at 16 coordinates (2.661 versus 2.845 for RSSPCA on the scaled MSE). At 50\%, RSSPCA's MSE was 4.162, 24.5\% below PCA and 18.5\% below its no-smoothing arm. Both stronger training corruption and higher dimensional budgets have visible effects on representation fidelity. The held-out input is already clean. These numbers measure whether the subspace fitted under corruption can represent clean faces.

For a fixed fit, the first $k$ columns are prefixes of the same loading matrix. Increasing $k$ enlarges the projection span, so least-squares reconstruction error can only decrease in exact arithmetic. Identification accuracy has no such monotonic guarantee because adding a score changes neighbor distances and votes.

\clearpage
\subsection{Errors, examples and external validity}

\reportfigure{face_reconstruction}{Per-image reconstruction MSE multiplied by $10^3$ for the three PCA-family arms at 16 coordinates, measured before clipping. Fits use 50\% corrupted training with seed 8301. Each example is the smallest held-out source index for one of the five identities fixed before scoring. Each marker is one error on one clean input. Table~\ref{tab:facemse} summarizes all 80 held-out views.}{fig:facerecon}

At seed 8301 with 50\% training corruption, RSSPCA identified 50 of 80 faces correctly. The two held-out views of identity 12 were assigned to 3, both views of 22 to 8, and both views of 38 to 28. Across five severe masks, seven identities had both test views correct in every fit, while one identity had neither view correct in any fit. Each identity contributes only two distinct images, so these outcomes describe the observed split. They do not isolate which visual feature drove a particular error.

\textbf{Limitations.} The cropped, frontal data and $16\times16$ reduction restrict the range of acquisition conditions. The shared patch treatment changes the fitted directions and the neighbor exemplars together. The experiment cannot attribute a recognition error to just one pathway. Among 880 primary-budget RSSPCA classifications, 220 used the fixed class-vote tie rule. Those counts include repeated predictions of the same 80 images. Higher-resolution scaling, occluded test views and sensitivity to other tie rules remain unmeasured.

\clearpage
\subsection{Relation to the published experiment and digit study}

Wang et al.\ \citep[Secs.~4.1--4.2]{wang2025} resize ORL faces to $56\times46$. Their identification study uses 7/3 images per identity, training-derived feature z-scores, 30 directions and nearest-neighbor classification across three repetitions. Our $16\times16$ Olivetti setup uses a fixed 6/2/2 split, a final 8/2 refit, and mean-centered inputs with primary budget 16 and shared 3-NN. For reconstruction, the source fits all 400 ORL images after perturbing 80 and measures errors on the 320 uncorrupted images in that pool. Our errors use held-out inputs and targets. These protocol differences make the present results an adaptation with a separate interpretation.

\begin{table}[H]\small\centering
\caption{Secondary transfer evidence from Prompt 01. Optdigits accuracy (\%) at eight coordinates, with a writer-disjoint holdout of 1,797 images. These metrics use a different dataset and protocol.}
\label{tab:digitsecondary}
\begin{tabular}{@{}lrrr@{}}
\toprule
Method & Clean & 25\% & 50\%\\
\midrule
PCA & 94.55 & 92.95 & 91.13\\
Fisher LDA & 95.05 & 93.98 & 91.89\\
No smoothness & 94.60 & 92.92 & 91.09\\
RSSPCA & 94.77 & 93.12 & 90.76\\
\bottomrule
\end{tabular}
\end{table}

That experiment predates the face request. It tuned on the official UCI Optdigits development set \citep{uci1998} and tested different writers, whereas the face test uses different images of known identities. Its advantage over PCA at 25\% and deficit at 50\% show that a favorable face result cannot be transferred to another dataset by assumption. The original digits protocol, all 165 prediction arrays and their verification are preserved beside the archived digit report.

\subsection*{Conclusion}
The 2025 RSSPCA method was numerically feasible on $16\times16$ faces. Under severe training corruption its graph term reduced loading roughness, and its fitted span represented clean test faces more accurately than the matched PCA-family controls. Recognition gains were smaller, varied across masks, and remained below Fisher LDA. The face and digit results jointly show why structural effects, reconstruction and class decisions should be measured separately.

\clearpage
\section{What I Have Learnt from This AI Assignment}

Compared with my previous AI assignments, I found that this assignment was not just about getting Codex to produce a working project. I also had to think more carefully about whether the project itself was a good choice. At first, I wondered whether choosing a newer and more complicated algorithm would make the assignment better. However, after looking into different PCA/LDA methods, I realized that difficulty and publication date are not everything. An algorithm might look impressive on paper, but if it is too difficult to implement properly or does not allow meaningful experiments, it might not be a suitable choice. I learned that choosing a research topic involves finding a balance between technical difficulty, relevance, feasibility, and what can actually be demonstrated.

One decision I spent quite some time thinking about was the dataset. Initially, Codex implemented Robust Sparse Smooth PCA (RSSPCA) for handwritten-digit recognition. The implementation worked, and the topic met the assignment requirements. However, I noticed that it was quite similar to my previous assignment, which also used handwritten digits. Even though the algorithms were completely different, I felt that using a similar dataset and producing similar types of graphs would make the two projects look repetitive. I considered whether I should start everything again with a different topic. Eventually, I decided to keep RSSPCA and change the main application to face recognition, which was also closer to the original research paper. This experience taught me that the first project design is not necessarily the best one, and sometimes changing direction does not mean throwing away all the work already done.

I also gained a better understanding of PCA and LDA beyond what I had learned in class. Previously, I mainly understood PCA as a way of reducing the number of features while retaining as much information as possible. Through this assignment, I learned how RSSPCA modifies the traditional approach by considering robustness, sparsity, and spatial smoothness. I also realized that understanding the mathematical idea is only part of the challenge. Actually implementing the algorithm involves many details, such as numerical stability, iterative optimization, stopping conditions, and computational cost. This became more obvious when moving from small digit images to face images, where the number of features could be much larger. It made me appreciate why checking whether an algorithm is practical should come before running a large experiment.

Another thing I found interesting was that a more advanced algorithm does not always give better results. In the initial digit experiments, RSSPCA produced smoother feature patterns, but its recognition accuracy was not consistently better than ordinary PCA or LDA. At first, I expected the newer method to perform better, but the results showed that this is not necessarily the case. I learned that an experiment should not be judged only by whether the proposed method wins. Understanding why different methods behave differently, what trade-offs they introduce, and under what conditions they work well is just as important. I also paid more attention to the graphs in this assignment. Instead of simply making them look professional, I wanted each figure to explain something useful about the algorithm or the experimental results.

Finally, this assignment helped me improve how I work with AI coding agents. From my previous assignments, I had already learned the importance of \texttt{AGENTS.md}, detailed prompts, and specialized skills. This time, I realized that even with well-written instructions, I should not simply accept every decision Codex makes. I still need to question whether its choices make sense, whether the experiments are fair, and whether the report actually demonstrates what the assignment asks for. Looking back, I think the biggest lesson for me was that using AI effectively is not just about telling it what to do. It is also about knowing when to question its decisions, when to improve the plan, and how to judge whether the final work is convincing.

\clearpage
\section{Webpage Link of the Source Codes}
\noindent\colorbox{pale}{\parbox{\dimexpr\linewidth-2\fboxsep\relax}{%
\textbf{Public source repository}\par
\url{https://github.com/Nvb-flipped/cisc3024-rsspca-face-recognition}\par
\textbf{GitHub Wiki}\par
\url{https://github.com/Nvb-flipped/cisc3024-rsspca-face-recognition/wiki}}}\par

\begin{table}[H]\small
\caption{Evidence available to a reviewer. Paths start at the project root.}
\begin{tabularx}{\linewidth}{@{}P{56mm}Y@{}}
\toprule
Path & Evidence\\
\midrule
\path{docs/face_pivot_decision.md}, \path{docs/face_study_design.md} &
Human request, source attribution, training-only gate, frozen study decisions and dated execution record.\\
\path{configs/faces.json}, \path{src/face_study.py} &
Fixed split, mask and tuning settings with distinct prepare and held-out phases.\\
\path{results/faces/} &
Split and selection, frozen weights and masks, predictions, reconstruction errors, summaries and replay audit.\\
\path{report/optdigits/}, \path{configs/study.json} &
Archived digit PDFs and source, original settings, and links to the unchanged digit evidence.\\
\path{src/face_summary.py}, \path{src/face_figures.py} &
Recreate the face tables and vector figures from the archived records.\\
\bottomrule
\end{tabularx}
\end{table}

The face replay uses \texttt{python -m src.face\_study} with stages \texttt{fetch}, \texttt{tune}, \texttt{prepare}, \texttt{final} and \texttt{audit}. It can compare a fresh output directory to the saved face results. The raw Olivetti archive is obtained separately and checked by SHA-256. The public figures report measured errors, generated masks and learned coefficients. The README gives the commands to reproduce the study and regenerate the local illustration galleries from the attributed source.

\clearpage
\renewcommand{\bibsection}{\section*{References}}
\bibliographystyle{plainnat}
\begingroup\raggedright
\bibliography{references}
\endgroup
\end{document}
